\clearpage\thispagestyle{empty}
\pdfbookmark[0]{Quellen zu Gauß’ Besprechung von Laplace}{laplace1808-sources}
\begin{center}\textbf{Quellen zu Gauß’ Besprechung von Laplace}\end{center}
\noindent Diese KI-Übertragung ins heutige Deutsch folgt Gauß’ Besprechung auf S.~500--503 der \href{https://gdz.sub.uni-goettingen.de/id/PPN236007467}{\textit{Werke}, Bd.~VI} (Göttingen, 1874). Sie erschien ohne Namensnennung unter der Überschrift \emph{Paris} in den \textit{Göttingischen gelehrten Anzeigen}, Stück~119, 23. Juli 1808, S.~\href{https://api.digitale-sammlungen.de/iiif/image/v2/bsb10538674_00155/full/1800,/0/default.jpg}{1185}--\href{https://api.digitale-sammlungen.de/iiif/image/v2/bsb10538674_00162/full/1800,/0/default.jpg}{1192}. Das verlinkte Exemplar stammt aus der Bayerischen Staatsbibliothek. Einige Kolumnentitel tragen das Datum 25. Juli.
\par\smallskip\noindent Gauß bespricht die dritte Ausgabe von Pierre-Simon Laplaces \href{https://www.loc.gov/item/17002967/}{\textit{Exposition du système du monde}} (Paris, Courcier, 1808). Der Titel bedeutet \emph{Darstellung des Weltsystems}. Zwei Zeichen in der Tabelle der säkularen Änderungen sehen in der Werkausgabe wie ein kleines c aus. Im Eintrag zur großen Halbachse des Jupiter steht zudem ein Zwischenraum. Der \href{https://api.digitale-sammlungen.de/iiif/image/v2/bsb10538674_00159/full/1800,/0/default.jpg}{Erstdruck} zeigt an diesen Stellen Nullen und ein Dezimalkomma. Die übertragenen Tabellen folgen der Werkausgabe.
\par\smallskip\noindent Die französischen Worte, die Gauß im Absatz zum fünften Buch aus Laplace zitiert, lauten \emph{Le tableau des progrès de la plus sublime des sciences naturelles, toujours croissans au milieu même des révolutions des empires, pourra consoler des malheurs dont les récits remplissent les annales de tous les peuples.} KI-Übertragung und KI-Satz.
\clearpage
